
\section{AT-03: Hydrogenic Stark shift from closure}
\source{2026 Atomic Assignment 3, Question 3(a)--(c).}
\subsection*{Step 1: write the perturbation and parity argument}
An atom in a uniform electric field \(\mathcal E\hat z\)
has perturbation \(H_E=e\mathcal E z\) for an electron,
because its dipole is \(-e\mathbf r\).
A non-degenerate \(1s\) state has even parity, so
\(\langle1s|z|1s\rangle=0\) and the first-order Stark
shift vanishes. In second order,
\[
\Delta E_{1s}^{(2)}=
\mathcal E^2\sum_{m\ne1s}
\frac{|\langle m|ez|1s\rangle|^2}{E_{1s}-E_m}
=-\frac12\alpha\mathcal E^2.
\]
Hence, in atomic units,
\[
\boxed{\alpha=2\sum_{m\ne1s}
\frac{|\langle m|z|1s\rangle|^2}{E_m-E_{1s}}}.
\]

\subsection*{Step 2: approximate all denominators}
As requested in the assignment, take
\(E_m-E_1\simeq E_2-E_1=3Z^2/8\) Hartree.
Using completeness,
\[
\sum_{m\ne1s}|\langle m|z|1s\rangle|^2
=\langle1s|z^2|1s\rangle
-|\langle1s|z|1s\rangle|^2 .
\]
The second term is zero by parity. Spherical symmetry
gives \(\langle z^2\rangle=\langle r^2\rangle/3\).
For the hydrogenic \(1s\) orbital,
\(\langle r^2\rangle=3/Z^2\) in Bohr-radius units.
Therefore \(\langle z^2\rangle=1/Z^2\), and
\[
\boxed{\alpha_{\rm approx}
=\frac{2(1/Z^2)}{3Z^2/8}
=\frac{16}{3Z^4}\ a_0^3}.
\]

\subsection*{Step 3: numerical checks and interpretation}
For hydrogen, \(Z=1\):
\[
\alpha_{\rm approx}=\frac{16}{3}a_0^3
=5.333\,a_0^3\simeq7.903\times10^{-25}\ {\rm cm}^3,
\]
using \(a_0^3=1.48185\times10^{-25}\ {\rm cm}^3\).
The exact \(1s\) result is
\(\alpha_{\rm exact}=4.5\,a_0^3\), so this crude
denominator approximation overestimates it by \(18.5\%\).

For hydrogen-like \(\mathrm{Rb}^{36+}\) (\(Z=37\)),
\[
\boxed{\alpha_{\rm approx}=
\frac{16}{3(37)^4}a_0^3
=2.846\times10^{-6}a_0^3
\simeq4.217\times10^{-31}\ {\rm cm}^3}.
\]
The tightly bound single \(1s\) electron is much less
polarisable than valence electrons in \emph{neutral} Rb.
\checkmethod{Use parity before a long sum; in closure,
exclude the initial state explicitly. The \(Z^{-4}\)
scaling is a powerful plausibility test.}
